\Needspace{10\baselineskip}
\centerline{\textsc{Ceres und Pallas.}}

\medskip
Monatliche Correspondenz. Band IX, S.\ 246—250. März 1804.

\medskip
Im Oktoberheft des vergangenen Jahres haben wir auf S.\ 370 bereits eine Ephemeride des geozentrischen Laufs der \emph{Ceres} für das Jahr 1804 veröffentlicht, die Dr.\ \textsc{Gauß} aus seinen zum neunten Mal verbesserten Bahnelementen berechnet hatte. Im Dezemberheft haben wir auf S.\ 535 eine Karte dieses Laufs von Inspektor \textsc{Harding} in Lilienthal angekündigt. Sie liegt nun diesem Heft bei.

{\rightskip=0pt plus 1em
Dr.\ \textsc{Gauß} hat seither die Elemente der \emph{Pallas}-Bahn anhand der neuesten Beobachtungen noch einmal korrigiert und aus diesen zum siebten Mal verbesserten Bahnelementen die folgende Ephemeride berechnet. Auch auf diese werden wir eine von Inspektor \textsc{Harding} entworfene Karte des geozentrischen Laufs dieses Planeten folgen lassen. [Die Ephemeride reicht vom 30. April 1804 bis zum 19. Januar 1805 und gibt in der bisher verwendeten Form die Rektaszension, die Deklination, die Entfernung von der Erde und die Lichtstärke an.]\par}


An einigen Stellen ist der Unterschied zur Ephemeride beträchtlich, die Prof.\ \textsc{Bode} im Berliner Astronomischen Jahrbuch für 1806 auf S.\ 91 angegeben hat. \emph{Pallas} wird sich in diesem Jahr wieder etwas leichter beobachten lassen als im vergangenen. Ihre größte Lichtstärke vom 20. Juni des vergangenen Jahres erreicht sie dieses Jahr am 23. Juli, und am 7. September wird sie noch um mehr als $\frac{1}{4}$ heller. Dazu kommen die völlige Dunkelheit der Nächte und der Umstand, dass \emph{Pegasus} und \emph{Wassermann} weniger sternreich sind als die Gegend, in der sich der Planet im vergangenen Jahr aufhielt. Die Lichtstärke, bei der er im vergangenen Jahr zum ersten Mal gefunden wurde, erreichte er diesmal schon am 5. Juni, und am 18. Mai erreichte er die Lichtstärke vom 10. Oktober 1803, als Dr.\ \textsc{Olbers} ihn zum letzten Mal beobachtete.

Es folgen die neuesten Elemente (VII) der \emph{Pallas}-Bahn, aus denen Dr.\ \textsc{Gauß} die obige Ephemeride berechnet hat.

\begin{center}
\begingroup\fontsize{9}{11}\selectfont\setlength{\tabcolsep}{3pt}
\begin{tabular}{lr}
Epoche (Seeberg) 1803 . . . . . . & $221^\circ\ 29'\ 32''0$ \\
tägliche tropische Bewegung . . . & $770''0446$ \\
jährliche tropische Bewegung . . . & $78^\circ\ 4'\ 26''3$ \\
Aphel 1803 . . . . . . . & $301\ 17\ 34.4$ \\
Exzentrizität . . . . . . . . . & $0.2457396$ \\
Logarithmus der großen Halbachse  & $0.4423790$ \\
Aufsteigender Knoten 1803 . . . . . . & $172^\circ\ 28'\ 13''7$ \\
Neigung der Bahn . . . . . . . & $34\ 38\ 1.1$ \\
\end{tabular}\endgroup
\end{center}
\SourceLeaf{252}%

Im letzten Dezemberheft haben wir auf S.\ 536 aus dem großen Sternkatalog von \textsc{Piazzi} ein kleines Verzeichnis der Sterne veröffentlicht, die auf dem Weg der \emph{Ceres} oder in seiner Nähe liegen werden. Ebenso geben wir hier eine kleine Auswahl der Sterne aus diesem Katalog an, die der \emph{Pallas} sehr nahe kommen und mit ihr verglichen werden können.

\begin{flushright}
[\textsc{von Zach}.]
\end{flushright}

\bigskip
